\documentclass[11pt,aspectratio=169]{beamer}
\usetheme{SimplePlus}

\usepackage[T1]{fontenc}
\usepackage[utf8]{inputenc}
\usepackage{amsmath,amssymb,mathtools}
\usepackage{booktabs}
\usepackage{array}
\usepackage{appendixnumberbeamer}
\usepackage{hyperref}

\newcommand{\E}{\mathbb{E}}
\newcommand{\Var}{\mathrm{Var}}
\newcommand{\Cov}{\mathrm{Cov}}
\newcommand{\Prob}{\mathbb{P}}
\newcommand{\ATE}{\mathrm{ATE}}
\newcommand{\LATE}{\mathrm{LATE}}
\newcommand{\indep}{\perp\!\!\!\perp}
\newcommand{\derivbutton}[1]{\vspace{0.4em}\hfill\hyperlink{app:#1}{\beamerbutton{Appendix details}}}
\newcommand{\backbutton}[1]{\vspace{0.4em}\hfill\hyperlink{main:#1}{\beamerbutton{Back to main slide}}}

\title{Applied Econometrics: Session 9}
\subtitle{Regression Discontinuity Design}
\author{PhD Intensive Course}
\institute{Lecture 9}
\date{\today}

\begin{document}

\begin{frame}
  \titlepage
\end{frame}

\begin{frame}{How This Lecture Fits the Course}
  Earlier lectures used controls, instruments, matching, DiD, fixed effects, and synthetic controls to construct credible counterfactuals.
  \vspace{0.4em}

  Regression discontinuity uses a different source of credibility:
  \begin{block}{Core idea}
    If treatment changes discontinuously at a cutoff, units just below and just above the cutoff can be compared locally.
  \end{block}

  \begin{itemize}
    \item The running variable is observed and determines eligibility.
    \item The cutoff creates a jump in treatment probability.
    \item The causal estimate is local: it applies at the cutoff, not automatically far away.
  \end{itemize}
\end{frame}

\begin{frame}{Roadmap for Today}
  \tableofcontents
\end{frame}

\section{Cutoff Intuition}

\begin{frame}{A Threshold as a Natural Experiment}
  Examples of cutoff rules:
  \begin{itemize}
    \item students pass an exam if score \(\geq 50\);
    \item firms receive a subsidy if employment is below 250 workers;
    \item people become legally eligible for a program at a certain age;
    \item municipalities receive extra funding if population is below a threshold.
  \end{itemize}
  \begin{block}{Why this may help}
    A person who scores 49.9 and a person who scores 50.1 are probably very similar, except that the rule treats them differently.
  \end{block}
\end{frame}

\begin{frame}{The Running Variable}
  Let \(R_i\) be the running variable and \(c\) the cutoff.
  \[
    X_i \equiv R_i-c
  \]
  is the centered running variable.
  \begin{itemize}
    \item \(X_i<0\): unit is below the cutoff.
    \item \(X_i=0\): unit is exactly at the cutoff.
    \item \(X_i>0\): unit is above the cutoff.
  \end{itemize}
  \begin{block}{Practical habit}
    Center the running variable before estimating RD. Then the intercepts of local regressions are values at the cutoff.
  \end{block}
\end{frame}

\begin{frame}{Sharp RD Treatment Rule}
  \hypertarget{main:sharp-rule}{}
  A sharp RD has deterministic assignment:
  \[
    D_i = \mathbf{1}\{R_i \geq c\}
    =
    \mathbf{1}\{X_i \geq 0\}.
  \]
  \begin{itemize}
    \item Everyone below the cutoff is untreated.
    \item Everyone above the cutoff is treated.
    \item The treatment probability jumps from 0 to 1 at \(c\).
  \end{itemize}
  \begin{block}{Important}
    Sharp means the rule fully determines treatment. It does not mean the design is automatically valid.
  \end{block}
\end{frame}

\begin{frame}{Potential Outcomes Near a Cutoff}
  Each unit has two potential outcomes:
  \[
    Y_i(1), \qquad Y_i(0).
  \]
  We observe
  \[
    Y_i = D_iY_i(1)+(1-D_i)Y_i(0).
  \]
  \begin{itemize}
    \item Above the cutoff, sharp RD reveals treated outcomes.
    \item Below the cutoff, sharp RD reveals untreated outcomes.
    \item The missing object is the counterfactual outcome on the other side of the cutoff.
  \end{itemize}
\end{frame}

\begin{frame}{The Visual Logic}
  Imagine plotting average outcome against \(X_i=R_i-c\).
  \begin{itemize}
    \item Away from the cutoff, treated and untreated units may be very different.
    \item Very close to the cutoff, the running variable is nearly the same.
    \item A discontinuous jump in the outcome at \(X=0\) is evidence of a treatment effect.
  \end{itemize}
  \begin{block}{RD picture in words}
    Fit the smooth left-hand relationship up to the cutoff and the smooth right-hand relationship down to the cutoff. The vertical gap at zero is the RD estimate.
  \end{block}
\end{frame}

\begin{frame}{What RD Identifies}
  RD is not mainly about comparing all treated units with all untreated units.
  \[
    \text{all above cutoff} \quad \text{vs.} \quad \text{all below cutoff}
  \]
  is usually a bad comparison because high-\(R\) and low-\(R\) units differ.
  \vspace{0.4em}

  RD is about the limit comparison:
  \[
    R_i \downarrow c
    \quad \text{vs.} \quad
    R_i \uparrow c.
  \]
  \begin{block}{Local target}
    RD estimates the treatment effect for units at the cutoff, or for units close enough that the cutoff comparison is credible.
  \end{block}
\end{frame}

\begin{frame}{Why ``Barely Passed'' Is Not ``Passed by a Lot''}
  Suppose an exam score determines graduation.
  \begin{itemize}
    \item Score 90 vs. score 40: very different students.
    \item Score 50.1 vs. score 49.9: likely similar students.
    \item The cutoff comparison removes much of the concern that ability drives the result.
  \end{itemize}
  \begin{block}{Student warning}
    RD is credible because it is narrow. Do not turn a local design into a broad claim without extra evidence.
  \end{block}
\end{frame}

\section{Sharp RD Identification}

\begin{frame}{The Continuity Assumption}
  \hypertarget{main:continuity}{}
  For \(d\in\{0,1\}\), define
  \[
    \mu_d(r)=\E[Y_i(d)\mid R_i=r].
  \]
  The key identifying assumption is:
  \[
    \lim_{r\to c^-}\mu_d(r)
    =
    \lim_{r\to c^+}\mu_d(r),
    \qquad d=0,1.
  \]
  \begin{itemize}
    \item Potential outcomes may slope with the running variable.
    \item They must not jump at the cutoff for reasons other than treatment.
    \item The observed outcome can jump because treatment jumps.
  \end{itemize}
  \derivbutton{sharp-id}
\end{frame}

\begin{frame}{Sharp RD Estimand}
  \hypertarget{main:sharp-estimand}{}
  Under continuity and sharp assignment:
  \[
    \tau_{\mathrm{SRD}}
    =
    \lim_{r\to c^+}\E[Y_i\mid R_i=r]
    -
    \lim_{r\to c^-}\E[Y_i\mid R_i=r].
  \]
  This equals
  \[
    \E[Y_i(1)-Y_i(0)\mid R_i=c].
  \]
  \begin{block}{Interpretation}
    The estimand is a treatment effect at the cutoff, not necessarily the average treatment effect for the full population.
  \end{block}
  \derivbutton{sharp-id}
\end{frame}

\begin{frame}{Local Randomization Intuition}
  Near the cutoff, small differences in \(R_i\) can be as good as random.
  \begin{itemize}
    \item A student may barely pass or barely fail because of small test-day noise.
    \item A person may be just above or below an age threshold because of date of birth.
    \item A firm may just exceed an administrative size threshold.
  \end{itemize}
  \begin{block}{But be precise}
    Local randomization is an intuition. The formal RD assumption is continuity of potential outcomes at the cutoff.
  \end{block}
\end{frame}

\begin{frame}{What Can Break Continuity?}
  Continuity fails if something else jumps at the same cutoff.
  \begin{itemize}
    \item Another policy starts at exactly the same threshold.
    \item Measurement of the outcome changes at the cutoff.
    \item People sort around the threshold using private information.
    \item The cutoff itself is chosen because expected outcomes jump there.
  \end{itemize}
  \begin{block}{Design question}
    If treatment did not change at the cutoff, would the expected outcome have moved smoothly through it?
  \end{block}
\end{frame}

\begin{frame}{A Numeric Cutoff Example}
  Imagine a scholarship for scores \(\geq 70\).
  \[
  \begin{array}{lcc}
  \toprule
  & \text{Mean outcome just below 70} & \text{Mean outcome just above 70} \\
  \midrule
  \text{Observed graduation rate} & 0.62 & 0.69 \\
  \bottomrule
  \end{array}
  \]
  The sharp RD estimate is
  \[
    0.69-0.62=0.07.
  \]
  \begin{block}{Correct wording}
    The scholarship increased graduation by 7 percentage points for students at the score cutoff, if potential graduation rates would otherwise have been continuous at 70.
  \end{block}
\end{frame}

\begin{frame}{What Not to Estimate}
  A global mean comparison:
  \[
    \E[Y_i\mid R_i\geq c]-\E[Y_i\mid R_i<c]
  \]
  is generally not an RD estimand.
  \begin{itemize}
    \item Units far above and far below the cutoff can be very different.
    \item The running variable itself often strongly predicts the outcome.
    \item RD asks for a local discontinuity, not a global level difference.
  \end{itemize}
  \begin{block}{Exam phrase}
    ``Close to the cutoff'' is not a cosmetic detail. It is the source of identification.
  \end{block}
\end{frame}

\section{Local Linear Estimation}

\begin{frame}{Center First}
  Define
  \[
    X_i=R_i-c,
    \qquad
    Z_i=\mathbf{1}\{X_i\geq 0\}.
  \]
  Centering makes the cutoff equal to zero.
  \begin{itemize}
    \item Left of cutoff: \(Z_i=0\).
    \item Right of cutoff: \(Z_i=1\).
    \item At \(X_i=0\), slopes vanish and intercepts have a direct cutoff interpretation.
  \end{itemize}
\end{frame}

\begin{frame}{Local Linear RD Regression}
  \hypertarget{main:local-linear}{}
  A standard sharp RD specification is
  \[
    Y_i
    =
    \beta_0+\beta_1X_i+\tau Z_i+\beta_2(Z_iX_i)+u_i.
  \]
  \begin{itemize}
    \item Below cutoff: \(\E[Y_i\mid X_i]=\beta_0+\beta_1X_i\).
    \item Above cutoff: \(\E[Y_i\mid X_i]=(\beta_0+\tau)+(\beta_1+\beta_2)X_i\).
    \item At \(X_i=0\), the right-left jump is \(\tau\).
  \end{itemize}
  \begin{block}{Interpretation}
    This is like fitting two lines, one on each side, and comparing their intercepts at the cutoff.
  \end{block}
  \derivbutton{local-linear}
\end{frame}

\begin{frame}{Why Local, Not Global?}
  If the true relationship between \(Y\) and \(R\) is curved, a line fit far from the cutoff can distort the predicted value at the cutoff.
  \begin{itemize}
    \item A very wide window gives more observations but more functional-form bias.
    \item A very narrow window gives less bias but noisier estimates.
    \item RD estimation is therefore a bias-variance tradeoff.
  \end{itemize}
  \begin{block}{Practical rule}
    The chosen bandwidth should be reported, justified, and varied in robustness checks.
  \end{block}
\end{frame}

\begin{frame}{Bandwidth}
  Use observations with
  \[
    |X_i|\leq h.
  \]
  The bandwidth \(h\) controls how local the comparison is.
  \begin{itemize}
    \item Small \(h\): close comparison, high variance.
    \item Large \(h\): more precision, more extrapolation.
    \item Main estimates should not rely on one fragile bandwidth choice.
  \end{itemize}
  \begin{block}{Student habit}
    Always ask: how many observations are inside the bandwidth on each side?
  \end{block}
\end{frame}

\begin{frame}{Triangular Kernel Weights}
  A common local weighting rule is the triangular kernel:
  \[
    K_i(h)
    =
    \mathbf{1}\{|X_i|\leq h\}
    \left(1-\frac{|X_i|}{h}\right).
  \]
  \begin{itemize}
    \item Units outside the bandwidth get zero weight.
    \item Units near the cutoff get the largest weight.
    \item Units near the edge of the window get small weight.
  \end{itemize}
  \begin{block}{Why this matches RD logic}
    The closest observations are the most credible comparisons.
  \end{block}
\end{frame}

\begin{frame}{Weighted Local Linear Estimator}
  \hypertarget{main:wls}{}
  Estimate
  \[
    \min_{\beta_0,\beta_1,\tau,\beta_2}
    \sum_i
    K_i(h)
    \left[
      Y_i-\beta_0-\beta_1X_i-\tau Z_i-\beta_2Z_iX_i
    \right]^2.
  \]
  The coefficient on \(Z_i\) is the estimated jump at the cutoff.
  \begin{block}{Key point}
    The regression is only an interpolation device near the cutoff. The design, not the linear model, is doing the causal work.
  \end{block}
  \derivbutton{local-linear}
\end{frame}

\begin{frame}{Why Local Linear Is the Workhorse}
  Local linear RD is popular because:
  \begin{itemize}
    \item it allows different slopes on the two sides;
    \item it focuses on the boundary point \(X=0\);
    \item it has better boundary behavior than a simple local mean;
    \item it is easy to explain and graph.
  \end{itemize}
  \begin{block}{Caution}
    High-order global polynomials can create artificial wiggles. They are usually a poor substitute for local estimation.
  \end{block}
\end{frame}

\begin{frame}{Standard Errors}
  RD uncertainty comes from estimating two limits with finite data.
  \begin{itemize}
    \item Use heteroskedasticity-robust standard errors as a default.
    \item If data are clustered, cluster at the assignment or sampling level.
    \item If the running variable is discrete with few support points, inference needs extra care.
  \end{itemize}
  \begin{block}{Reporting}
    Report the point estimate, standard error, confidence interval, bandwidth, kernel, and observations on both sides.
  \end{block}
\end{frame}

\section{Diagnostics}

\begin{frame}{Graphical RD Plot}
  A useful RD graph has:
  \begin{itemize}
    \item binned means of \(Y_i\) against centered \(X_i\);
    \item a vertical line at the cutoff;
    \item separate fitted curves or local lines on each side;
    \item enough scale information to see the amount of data near zero.
  \end{itemize}
  \begin{block}{Purpose}
    The graph should show whether the estimated jump is visible and whether the functional form looks plausible near the cutoff.
  \end{block}
\end{frame}

\begin{frame}{Covariate Balance at the Cutoff}
  Pre-treatment covariates should not jump at the cutoff.
  \[
    \lim_{r\to c^+}\E[W_i\mid R_i=r]
    -
    \lim_{r\to c^-}\E[W_i\mid R_i=r]
    \approx 0
  \]
  for predetermined \(W_i\).
  \begin{itemize}
    \item Age, prior test scores, pre-policy outcomes, baseline demographics.
    \item A covariate jump suggests sorting, another policy, or measurement problems.
  \end{itemize}
  \begin{block}{Do not overstate}
    Balance tests cannot prove validity. They can reveal obvious failures.
  \end{block}
\end{frame}

\begin{frame}{Manipulation and Density Checks}
  \hypertarget{main:density}{}
  The density of the running variable should be smooth through the cutoff.
  \begin{itemize}
    \item If people can precisely manipulate \(R_i\), many may bunch just on the favorable side.
    \item Examples: reported income thresholds, exam retakes, administrative score rounding.
    \item A density jump weakens the ``barely below vs. barely above'' comparison.
  \end{itemize}
  \begin{block}{McCrary-style idea}
    Check whether the number of observations changes discontinuously at the cutoff.
  \end{block}
  \derivbutton{density-check}
\end{frame}

\begin{frame}{Placebo Checks}
  RD placebo checks ask whether jumps appear where they should not.
  \begin{itemize}
    \item Use fake cutoffs where treatment does not change.
    \item Use pre-treatment outcomes as placebo outcomes.
    \item Use predetermined covariates as placebo outcomes.
  \end{itemize}
  \begin{block}{Interpretation}
    A placebo failure does not tell you exactly what is wrong, but it tells you the cutoff is not behaving like a clean local experiment.
  \end{block}
\end{frame}

\begin{frame}{Donut RD}
  Sometimes observations extremely close to the cutoff are suspicious.
  \begin{itemize}
    \item Scores may be rounded or manually adjusted.
    \item People may manipulate just enough to cross the rule.
    \item Administrative discretion may concentrate at the threshold.
  \end{itemize}
  A donut RD drops a small interval around the cutoff:
  \[
    \delta < |X_i| \leq h.
  \]
  \begin{block}{Use}
    Donut RD is a sensitivity check, not a license to ignore manipulation.
  \end{block}
\end{frame}

\begin{frame}{RD Reporting Checklist}
  A credible RD write-up should state:
  \begin{itemize}
    \item running variable and cutoff;
    \item why the rule changes treatment at that cutoff;
    \item sharp or fuzzy design;
    \item local estimand and population near the cutoff;
    \item bandwidth, kernel, and polynomial order;
    \item RD plot, covariate balance, and density/manipulation evidence;
    \item robustness to bandwidths and reasonable specifications.
  \end{itemize}
\end{frame}

\section{Fuzzy RD}

\begin{frame}{When RD Is Fuzzy}
  A fuzzy RD occurs when treatment probability jumps at the cutoff but not from 0 to 1.
  \[
    \lim_{r\to c^+}\Prob(D_i=1\mid R_i=r)
    \neq
    \lim_{r\to c^-}\Prob(D_i=1\mid R_i=r)
  \]
  but both probabilities may be between 0 and 1.
  \begin{itemize}
    \item Some eligible people do not take treatment.
    \item Some ineligible people receive treatment anyway.
    \item The cutoff is an instrument for treatment receipt.
  \end{itemize}
\end{frame}

\begin{frame}{Examples of Fuzzy RD}
  \begin{itemize}
    \item Legal drinking age increases access to alcohol, but some underage people drink.
    \item Passing an exam increases diploma receipt, but some failing students later receive credentials.
    \item Program eligibility increases participation, but some eligible people do not enroll.
    \item A scholarship cutoff increases aid receipt, but take-up is incomplete.
  \end{itemize}
  \begin{block}{Translation}
    Sharp RD estimates the effect of eligibility if eligibility equals treatment. Fuzzy RD estimates the effect of treatment for cutoff compliers.
  \end{block}
\end{frame}

\begin{frame}{Cutoff as an Instrument}
  Define the cutoff indicator:
  \[
    Z_i=\mathbf{1}\{R_i\geq c\}.
  \]
  In fuzzy RD, \(Z_i\) is an instrument for actual treatment \(D_i\).
  \begin{itemize}
    \item First stage: \(Z_i\) changes treatment probability at \(c\).
    \item Reduced form: \(Z_i\) changes the outcome at \(c\).
    \item Wald ratio: reduced-form jump divided by first-stage jump.
  \end{itemize}
  \begin{block}{Connection to Lecture 5}
    Fuzzy RD is IV where the instrument is local cutoff eligibility.
  \end{block}
\end{frame}

\begin{frame}{Fuzzy RD Estimand}
  \hypertarget{main:fuzzy-estimand}{}
  The fuzzy RD estimand is
  \[
    \tau_{\mathrm{FRD}}
    =
    \frac{
      \lim_{r\to c^+}\E[Y_i\mid R_i=r]
      -
      \lim_{r\to c^-}\E[Y_i\mid R_i=r]
    }{
      \lim_{r\to c^+}\E[D_i\mid R_i=r]
      -
      \lim_{r\to c^-}\E[D_i\mid R_i=r]
    }.
  \]
  Under the usual IV assumptions locally at the cutoff, this is
  \[
    \E[Y_i(1)-Y_i(0)\mid \text{complier at }c].
  \]
  \derivbutton{fuzzy-wald}
\end{frame}

\begin{frame}{Fuzzy RD Assumptions}
  For a local Wald interpretation, we need:
  \begin{enumerate}
    \item continuity of potential outcomes at the cutoff;
    \item a nonzero first-stage jump in treatment;
    \item exclusion: crossing the cutoff affects \(Y\) only through \(D\);
    \item monotonicity: crossing the cutoff does not make anyone less likely to receive treatment.
  \end{enumerate}
  \begin{block}{Interpretation}
    The estimate is local twice: local to the cutoff and local to compliers whose treatment status is changed by crossing the cutoff.
  \end{block}
\end{frame}

\begin{frame}{Numeric Fuzzy RD Example}
  \[
  \begin{array}{lcc}
  \toprule
  & \text{Just below cutoff} & \text{Just above cutoff} \\
  \midrule
  \E[Y\mid R] & 50 & 56 \\
  \E[D\mid R] & 0.30 & 0.80 \\
  \bottomrule
  \end{array}
  \]
  Reduced-form jump:
  \[
    56-50=6.
  \]
  First-stage jump:
  \[
    0.80-0.30=0.50.
  \]
  Fuzzy RD estimate:
  \[
    \frac{6}{0.50}=12.
  \]
\end{frame}

\begin{frame}{Sharp vs. Fuzzy RD}
  \small
  \begin{tabular}{>{\raggedright\arraybackslash}p{0.25\linewidth}>{\raggedright\arraybackslash}p{0.30\linewidth}>{\raggedright\arraybackslash}p{0.30\linewidth}}
    \toprule
    Feature & Sharp RD & Fuzzy RD \\
    \midrule
    Treatment rule & \(D_i=\mathbf{1}\{R_i\geq c\}\) & cutoff shifts probability of \(D_i\) \\
    Jump used & outcome jump & outcome jump divided by treatment jump \\
    Interpretation & effect at cutoff & complier effect at cutoff \\
    Analogy & local experiment & local IV design \\
    Main danger & discontinuity from other causes & weak first stage or exclusion failure \\
    \bottomrule
  \end{tabular}
\end{frame}

\section{Simulation and Review}

\begin{frame}{Python Example for This Lecture}
  The companion script is:
  \[
    \texttt{lectures/code/lecture-9-rdd.py}.
  \]
  It simulates:
  \begin{itemize}
    \item sharp RD data with a known treatment effect at the cutoff;
    \item local linear estimates under several bandwidths;
    \item fuzzy RD with incomplete take-up and a local Wald estimate;
    \item covariate and density diagnostics for manipulation checks.
  \end{itemize}
  \begin{block}{Use in class}
    Ask students to change the bandwidth and curvature, then predict how bias and variance should move before running the script.
  \end{block}
\end{frame}

\begin{frame}{What the Simulation Should Teach}
  \begin{itemize}
    \item The global treated-control gap is not the RD estimate.
    \item Local linear estimates get closer to the true cutoff effect when the window is reasonable.
    \item Very narrow windows can be unstable because little data remain.
    \item In fuzzy RD, the outcome jump must be scaled by the treatment jump.
    \item Manipulation shows up as suspicious bunching in the running variable.
  \end{itemize}
\end{frame}

\begin{frame}{Common Mistakes}
  \small
  \begin{itemize}
    \item Calling any threshold regression an RD without showing treatment jumps at the cutoff.
    \item Comparing all observations above and below the cutoff.
    \item Forgetting to center the running variable before interpreting the intercept.
    \item Using high-order global polynomials as the main design.
    \item Ignoring density and covariate balance diagnostics.
    \item Treating fuzzy RD as sharp RD.
    \item Claiming an ATE for everyone when the design identifies a cutoff effect.
  \end{itemize}
\end{frame}

\begin{frame}{Main Takeaways}
  \begin{enumerate}
    \item RD uses a cutoff rule to construct a local comparison.
    \item Sharp RD identifies the outcome jump at the cutoff as the treatment effect under continuity.
    \item Local linear regression estimates the two one-sided limits.
    \item Bandwidth choice controls the bias-variance tradeoff.
    \item Diagnostics should check graphs, covariates, and density around the cutoff.
    \item Fuzzy RD is a local IV design: reduced form divided by first stage.
  \end{enumerate}
\end{frame}

\begin{frame}{Source and Further Reading}
  This lecture is based on ideas from:
  \begin{itemize}
    \item Matheus Facure Alves, \emph{Causal Inference for the Brave and True}, Chapter 16, "Regression Discontinuity Design".
    \item Angrist and Pischke, \emph{Mostly Harmless Econometrics}, chapter on regression discontinuity.
    \item Angrist and Pischke, \emph{Mastering 'Metrics}, chapter on regression discontinuity.
    \item Lee and Lemieux (2010), "Regression Discontinuity Designs in Economics".
  \end{itemize}
  \vspace{0.3em}
  Handbook URL:
  \begin{itemize}
    \item \url{https://matheusfacure.github.io/python-causality-handbook/16-Regression-Discontinuity-Design.html}
  \end{itemize}
\end{frame}

\begin{frame}{What Students Should Be Able to Do}
  After this lecture, students should be able to:
  \begin{itemize}
    \item identify the running variable, cutoff, and treatment rule;
    \item write the sharp RD estimand using one-sided limits;
    \item explain the continuity assumption in words;
    \item estimate and interpret a local linear RD coefficient;
    \item explain the role of bandwidth and kernel weights;
    \item distinguish sharp and fuzzy RD;
    \item describe graphical, covariate, and manipulation diagnostics.
  \end{itemize}
\end{frame}

\appendix

\begin{frame}{Appendix Map}
  \begin{itemize}
    \item A1: Sharp RD identification.
    \item A2: Local linear RD coefficient.
    \item A3: Bandwidth bias-variance logic.
    \item A4: Fuzzy RD as local Wald/IV.
    \item A5: Density manipulation check.
  \end{itemize}
\end{frame}

\begin{frame}{Appendix A1: Sharp RD Identification}
  \hypertarget{app:sharp-id}{}
  \small
  In a sharp RD,
  \[
    D_i=\mathbf{1}\{R_i\geq c\}.
  \]
  Therefore,
  \[
    \lim_{r\to c^+}\E[Y_i\mid R_i=r]
    =
    \lim_{r\to c^+}\E[Y_i(1)\mid R_i=r],
  \]
  and
  \[
    \lim_{r\to c^-}\E[Y_i\mid R_i=r]
    =
    \lim_{r\to c^-}\E[Y_i(0)\mid R_i=r].
  \]
  By continuity of \(\E[Y_i(d)\mid R_i=r]\) at \(c\),
  \[
    \lim_{r\to c^+}\E[Y_i(1)\mid R_i=r]=\E[Y_i(1)\mid R_i=c],
  \]
  \[
    \lim_{r\to c^-}\E[Y_i(0)\mid R_i=r]=\E[Y_i(0)\mid R_i=c].
  \]
  Subtracting gives
  \[
    \tau_{\mathrm{SRD}}
    =
    \E[Y_i(1)-Y_i(0)\mid R_i=c].
  \]
  \backbutton{sharp-estimand}
\end{frame}

\begin{frame}{Appendix A2: Local Linear RD Coefficient}
  \hypertarget{app:local-linear}{}
  With \(X_i=R_i-c\) and \(Z_i=\mathbf{1}\{X_i\geq 0\}\), estimate
  \[
    Y_i
    =
    \beta_0+\beta_1X_i+\tau Z_i+\beta_2Z_iX_i+u_i.
  \]
  For \(X_i<0\), \(Z_i=0\), so
  \[
    \E[Y_i\mid X_i=x,Z_i=0]=\beta_0+\beta_1x.
  \]
  Taking the left limit at zero gives
  \[
    \lim_{x\to 0^-}\E[Y_i\mid X_i=x]=\beta_0.
  \]
  For \(X_i\geq 0\), \(Z_i=1\), so
  \[
    \E[Y_i\mid X_i=x,Z_i=1]=(\beta_0+\tau)+(\beta_1+\beta_2)x.
  \]
  Taking the right limit at zero gives \(\beta_0+\tau\). The jump is \(\tau\).
  \backbutton{local-linear}
\end{frame}

\begin{frame}{Appendix A3: Bandwidth Bias-Variance Logic}
  \hypertarget{app:bw-logic}{}
  Let \(\widehat{\tau}(h)\) be the RD estimate using bandwidth \(h\).
  A useful mental decomposition is
  \[
    \widehat{\tau}(h)-\tau
    =
    \underbrace{\text{bias}(h)}_{\text{curvature/extrapolation}}
    +
    \underbrace{\text{noise}(h)}_{\text{sampling variation}}.
  \]
  \begin{itemize}
    \item As \(h\) increases, more observations enter, so variance usually falls.
    \item As \(h\) increases, observations farther from the cutoff influence the fit, so bias may rise.
    \item As \(h\) shrinks, bias usually falls, but the estimate becomes noisy.
  \end{itemize}
  \[
    \text{MSE}(h)=\text{bias}(h)^2+\Var(\widehat{\tau}(h)).
  \]
  \backbutton{wls}
\end{frame}

\begin{frame}{Appendix A4: Fuzzy RD as Local Wald}
  \hypertarget{app:fuzzy-wald}{}
  Let \(Z_i=\mathbf{1}\{R_i\geq c\}\). Define the local reduced-form jump:
  \[
    \Delta_Y
    =
    \lim_{r\to c^+}\E[Y_i\mid R_i=r]
    -
    \lim_{r\to c^-}\E[Y_i\mid R_i=r].
  \]
  Define the local first-stage jump:
  \[
    \Delta_D
    =
    \lim_{r\to c^+}\E[D_i\mid R_i=r]
    -
    \lim_{r\to c^-}\E[D_i\mid R_i=r].
  \]
  Under continuity, exclusion, relevance, and monotonicity,
  \[
    \Delta_Y
    =
    \E[Y_i(1)-Y_i(0)\mid D_i(1)>D_i(0),R_i=c]\Delta_D.
  \]
  Hence
  \[
    \frac{\Delta_Y}{\Delta_D}
    =
    \E[Y_i(1)-Y_i(0)\mid \text{complier at }c].
  \]
  \backbutton{fuzzy-estimand}
\end{frame}

\begin{frame}{Appendix A5: Density Manipulation Check}
  \hypertarget{app:density-check}{}
  Let \(f_R(r)\) be the density of the running variable.
  A manipulation check asks whether
  \[
    \lim_{r\to c^-}f_R(r)
    =
    \lim_{r\to c^+}f_R(r).
  \]
  If the density jumps at \(c\), units may be sorting around the threshold.
  \vspace{0.4em}

  A simple classroom approximation:
  \begin{enumerate}
    \item bin the running variable into narrow intervals;
    \item count observations in each bin;
    \item look for a discontinuous count jump at the cutoff.
  \end{enumerate}
  Formal McCrary-style tests use local smoothing of the log density on both sides.
  \backbutton{density}
\end{frame}

\end{document}
